% Experiments --- setup, datasets, metrics, harness. Numbers live in
% Results; this section fixes the protocol. The LLM-benchmark table is a
% honest placeholder pending US-069 (H100 eval), per the data-real rule.
\section{Experimental Setup}
\label{sec:experiments}

\subsection{Datasets}
\paragraph{PASTIS-R}~\cite{garnot2021pastis} is the primary benchmark:
2{,}433 multi-temporal Sentinel-2 patches ($128\times128$ px), 18 crop
classes, metropolitan France, with five spatially disjoint folds.
Unless stated otherwise, segmentation models are evaluated on a
held-out fold and ensemble combiners on out-of-fold predictions.

\paragraph{Self-built PASTIS-style datasets (Italy, Germany).}
To measure transfer with real data where no PASTIS-format benchmark exists,
we built two datasets by hand in the same format as PASTIS-R. \emph{Italy
(Tuscany/Pianura, 2018)}: $1{,}438$ patches of $128\times128$\,px and
$103{,}350$ labeled parcels. \emph{Germany (Lower Saxony, DE4, 2023)}:
$1{,}246$ patches and $53{,}460$ parcels over a 14-month window ($\approx156$k
labeled parcels combined). Both follow the same recipe: EuroCrops declared-crop
polygons as labels, the full Sentinel-2 time series per patch (24 dates for Italy, ~41 for Germany) with a pixel-wise crop mask, an AlphaEarth summary per
parcel via Google Earth Engine, and a split into five geographic blocks for
leak-free evaluation. Constructing these datasets is what makes the
France$\rightarrow$Italy and France$\rightarrow$Germany transfer measurable at
all; they are versioned with DVC.

\paragraph{Sen4AgriNet (Catalonia)}~\cite{sykas2022sen4agrinet}
(CC-BY-SA-4.0) is used for the France$\rightarrow$Catalonia transfer
study. A Catalonia subset is versioned with DVC
(\texttt{data/sen4agrinet.dvc}).

\paragraph{EuroCropsML}~\cite{reuss2025eurocropsml} (CC-BY-SA-4.0)
supports the few-shot $k$-shot curve; the raw subset is versioned with
DVC (\texttt{data/transfer/eurocropsml.dvc}).

\paragraph{AgroMind}~\cite{ruan2025agromind} is used only to evaluate
the conversational layer. It has no training split (about 28{,}482 QA
pairs); fine-tuning on it would induce leakage, so it is strictly an
evaluation set.

\subsection{Metrics}
Segmentation is reported with mean intersection over union (mIoU) and
macro-averaged F1 at the parcel level; we also report pixel accuracy
where relevant. Macro averaging is the primary objective because it
penalises errors on rare and common crops equally, matching the
business objective. Ensemble combiners are compared by out-of-fold
macro F1 and accuracy.

\subsection{Evaluation harness and reproducibility}
All segmentation runs are scored by a shared evaluation harness
(\texttt{ml/eval/}) that enforces spatially disjoint folds and aborts on
geographic overlap between training and evaluation sets. A label-remap
module (\texttt{ml/eval/class\_remap.py}) projects predictions onto the
harmonised macro label space for cross-region comparison. Every figure
reported in Section~\ref{sec:results} is anchored to a versioned
artefact (JSON/CSV/registry) in the repository, recorded in
\texttt{\% src:} comments in the source; datasets and weights are
tracked with DVC and runs with MLflow.

\subsection{Conversational-layer benchmark protocol}
\label{sec:experiments-llm}
We evaluate the conversational layer along two independent axes. The
first is an \emph{engineering} axis --- backend interchangeability ---
which is already measured: behind a fixed perceiver, an identical
grounded query is issued to the three serving backends actually deployed
(Gemini~3.5~Flash in the cloud, plus Qwen3~MoE and Qwen~3.6-VL
on-premises), and we record whether the answer is identical and its
end-to-end latency. The measured result is reported in
Section~\ref{sec:results} (Table~\ref{tab:backend-switch}). The second
is a \emph{quality} axis --- an AgroMind LLM-as-judge protocol scoring the
frozen reasoner's answers for accuracy and grounding --- which depends on
a separate evaluation run on the H100 GPU (US-069) and is still pending;
we report the protocol here and defer the numbers rather than fabricate
them. No synthetic AgroMind scores are reported anywhere in this paper.
